\subsection{IPs con respuesta selectiva y diagnóstico TCP}\label{sec:selectiva}
\subsubsection{SCAMPER}
IP 37.79.226.46, destino 37.79.226.46. Se observó con icmp y no se observó con udp en las rondas y ventanas TTL comparables registradas. Resultado de repetición: corroborated\_\allowbreak{}observation.

\begingroup
\fontsize{9}{11}\selectfont
\addtocounter{table}{-1}
\begin{longtable}{p{1.5cm}p{1.8cm}p{2.2cm}p{8.7cm}}
\toprule
\textbf{Método} & \textbf{Repetición} & \textbf{IP observada} & \textbf{Run / ID RIPE} \\ \midrule
\endfirsthead
\toprule
\textbf{Método} & \textbf{Repetición} & \textbf{IP observada} & \textbf{Run / ID RIPE} \\ \midrule
\endhead
icmp & 1 & sí & scamper\_\allowbreak{}20261027T151550\_\allowbreak{}1b31253e \\
udp & 1 & no & scamper\_\allowbreak{}20261027T151628\_\allowbreak{}0033ccea \\
icmp & 2 & sí & scamper\_\allowbreak{}20261027T151818\_\allowbreak{}5de843b6 \\
udp & 2 & no & scamper\_\allowbreak{}20261027T151906\_\allowbreak{}e77fd62c \\
\bottomrule
\end{longtable}
\endgroup
Hipótesis: comportamiento dependiente del protocolo, posible filtrado o política de respuesta. El protocolo con respuesta permite observar esa IP en este experimento. No se demostró una ACL concreta: otra ruta, balanceo o pérdida también pueden explicar la ausencia. Las pruebas corroboran el contraste observado, no una regla de firewall.

\subsubsection{RIPE}
IP 200.0.16.25, destino 200.14.48.134. Se observó con TCP y no se observó con ICMP en las rondas y ventanas TTL comparables registradas. Resultado de repetición: corroborated\_\allowbreak{}observation.

\begingroup
\fontsize{9}{11}\selectfont
\addtocounter{table}{-1}
\begin{longtable}{p{1.5cm}p{1.8cm}p{2.2cm}p{8.7cm}}
\toprule
\textbf{Método} & \textbf{Repetición} & \textbf{IP observada} & \textbf{Run / ID RIPE} \\ \midrule
\endfirsthead
\toprule
\textbf{Método} & \textbf{Repetición} & \textbf{IP observada} & \textbf{Run / ID RIPE} \\ \midrule
\endhead
TCP & 1 & sí & ripe\_\allowbreak{}20261027T021029\_\allowbreak{}21cd8259 / 228007771 \\
ICMP & 1 & no & ripe\_\allowbreak{}20261027T021245\_\allowbreak{}9bda16d8 / 228009257 \\
TCP & 2 & sí & ripe\_\allowbreak{}20261027T021359\_\allowbreak{}987fefb9 / 228009551 \\
ICMP & 2 & no & ripe\_\allowbreak{}20261027T021549\_\allowbreak{}02ab18c4 / 228009740 \\
\bottomrule
\end{longtable}
\endgroup
Hipótesis: comportamiento dependiente del protocolo, posible filtrado o política de respuesta. El protocolo con respuesta permite observar esa IP en este experimento. No se demostró una ACL concreta: otra ruta, balanceo o pérdida también pueden explicar la ausencia. Las pruebas corroboran el contraste observado, no una regla de firewall.

\subsubsection{Discrepancia de recepción TCP}
En 36 ejecuciones Scamper TCP se correlacionaron respuestas ICMP por direcciones, IPID, puertos y tiempo, pero no se incorporaron como saltos aceptados en WARTS. Se conserva el resultado de la herramienta separado de la captura. El control revisó 1176 tramas TCP exteriores completas, incluidos RST de control, sin checksums TCP inválidos. Este conteo no equivale a sondas emitidas.

El control adicional con ICMP-rxerr finalizó con stop\_\allowbreak{}reason ERROR, mensaje de tipo no soportado y cero sondas. Se conserva como diagnóstico fallido y no cuenta como ruta válida. Su existencia no prueba filtrado ni corrige retrospectivamente los WARTS anteriores.

\subsubsection{Lectura de la cabecera citada en ICMP}
Ejemplo trazable: scamper\_\allowbreak{}20261026T204619\_\allowbreak{}270c47c3, trama 4. La respuesta exterior procede de 192.168.100.1 y tiene tipo/código ICMP 11/0. Dentro de ICMP se cita la solicitud original: destino 138.204.230.242, TTL 1, IPID 0x0001, UDP 32769 a 33435. El TTL exterior pertenece a la respuesta; no sustituye al TTL citado. La coincidencia se comprueba contra el paquete emitido, sus puertos, identificador y tiempo.

En las fichas, ip.proto 1 identifica ICMP, 6 TCP y 17 UDP; ICMP 8/0 es Echo Request. Los flags TCP 0x0002 y 0x0010 corresponden respectivamente a SYN y ACK. Las secuencias se leen con valores absolutos, evitando confundir la numeración relativa de la interfaz gráfica con el valor transmitido.

